\documentclass[11pt,a4paper]{article}

\usepackage[utf8]{inputenc}
\usepackage[T1]{fontenc}
\usepackage[french]{babel}
\usepackage[a4paper,margin=2.2cm]{geometry}
\usepackage{amsmath,amssymb,amsthm,mathtools}
\usepackage{lmodern}
\usepackage{enumitem}
\usepackage{hyperref}
\usepackage{xcolor}
\usepackage{fancyhdr}
\usepackage{stmaryrd}
\usepackage{mathrsfs}
\usepackage{array}
\usepackage{bm}

\hypersetup{
    colorlinks=true,
    linkcolor=black,
    urlcolor=blue,
    citecolor=black
}

\setlength{\parindent}{0pt}
\setlength{\parskip}{0.6em}

\pagestyle{fancy}
\fancyhf{}
\fancyhead[L]{Réduction --- Chapitre 17}
\fancyhead[R]{Trigonalisation et réduction plus fine}
\fancyfoot[C]{\thepage}

% Environnements
\newtheorem{definition}{Définition}[section]
\newtheorem{theoreme}[definition]{Théorème}
\newtheorem{proposition}[definition]{Proposition}
\newtheorem{corollaire}[definition]{Corollaire}
\newtheorem{lemme}[definition]{Lemme}
\theoremstyle{remark}
\newtheorem{remarque}[definition]{Remarque}
\newtheorem{exemple}[definition]{Exemple}
\newtheorem*{methode}{Méthode}
\newtheorem*{idee}{Idée}

% Commandes
\newcommand{\R}{\mathbb{R}}
\newcommand{\C}{\mathbb{C}}
\newcommand{\Q}{\mathbb{Q}}
\newcommand{\Z}{\mathbb{Z}}
\newcommand{\N}{\mathbb{N}}
\newcommand{\K}{\mathbb{K}}
\newcommand{\Vect}{\mathrm{Vect}}
\newcommand{\Ker}{\mathrm{Ker}}
\newcommand{\Img}{\mathrm{Im}}
\newcommand{\rg}{\mathrm{rg}}
\newcommand{\id}{\mathrm{Id}}
\newcommand{\Sp}{\mathrm{Sp}}
\newcommand{\tr}{\mathrm{tr}}
\newcommand{\diag}{\mathrm{diag}}
\newcommand{\Mnn}{M_n(\K)}
\newcommand{\GLn}{GL_n(\K)}
\newcommand{\chiA}{\chi_A}
\newcommand{\chil}{\chi_u}

\begin{document}

\begin{center}
    {\LARGE \textbf{Chapitre 17 --- Trigonalisation et réduction plus fine}}\\[0.4em]
    {\large Comprendre les matrices non diagonalisables}
\end{center}

\hrule
\vspace{1em}

\tableofcontents

\newpage

\section{Pourquoi la trigonalisation ?}

La diagonalisation est la situation idéale : on trouve une base de vecteurs propres et la matrice devient diagonale.  
Mais cette situation n’est pas toujours possible.

Par exemple, la matrice
\[
A=
\begin{pmatrix}
1 & 1\\
0 & 1
\end{pmatrix}
\]
n’est pas diagonalisable, bien qu’elle soit très simple.

La question naturelle devient alors :

\begin{center}
\textit{si une matrice n’est pas diagonalisable, peut-on quand même la réduire à une forme simple ?}
\end{center}

La réponse est oui, dans de nombreux cas : on cherche alors une base dans laquelle la matrice devient \textbf{triangulaire supérieure}.

\begin{idee}
La trigonalisation est une réduction plus faible que la diagonalisation, mais elle conserve déjà beaucoup d’informations utiles.
\end{idee}

Elle permet :
\begin{itemize}
    \item de lire immédiatement les valeurs propres sur la diagonale ;
    \item d’étudier les polynômes d’une matrice ;
    \item de comprendre les cas non diagonalisables ;
    \item de préparer des réductions plus fines.
\end{itemize}

\section{Matrices triangulaires}

\subsection{Définition}

\begin{definition}
Une matrice \(T=(t_{ij})\in M_n(\K)\) est dite \textbf{triangulaire supérieure} si
\[
i>j \Longrightarrow t_{ij}=0.
\]
Autrement dit, tous les coefficients situés strictement sous la diagonale sont nuls.
\end{definition}

\begin{exemple}
\[
\begin{pmatrix}
1 & 2 & -1\\
0 & 3 & 4\\
0 & 0 & 5
\end{pmatrix}
\]
est triangulaire supérieure.
\end{exemple}

\begin{definition}
Une matrice est dite \textbf{triangulaire inférieure} si tous les coefficients strictement au-dessus de la diagonale sont nuls.
\end{definition}

\subsection{Propriétés immédiates}

\begin{proposition}
Les valeurs propres d’une matrice triangulaire supérieure sont ses coefficients diagonaux.
\end{proposition}

\begin{proof}
Si
\[
T=(t_{ij})
\]
est triangulaire supérieure, alors
\[
XI_n-T
\]
est encore triangulaire supérieure, de diagonale
\[
X-t_{11},\dots,X-t_{nn}.
\]
Donc
\[
\chi_T(X)=\det(XI_n-T)=\prod_{i=1}^n (X-t_{ii}).
\]
Les racines du polynôme caractéristique sont donc exactement les coefficients diagonaux.
\end{proof}

\begin{proposition}
Le déterminant d’une matrice triangulaire supérieure est le produit de ses coefficients diagonaux :
\[
\det(T)=t_{11}\cdots t_{nn}.
\]
\end{proposition}

\begin{proof}
C’est le cas particulier du déterminant d’une matrice triangulaire.
\end{proof}

\begin{remarque}
Une matrice triangulaire est inversible si et seulement si tous ses coefficients diagonaux sont non nuls.
\end{remarque}

\section{Définition de la trigonalisation}

\subsection{Endomorphisme trigonalisable}

\begin{definition}
Soit \(E\) un \(\K\)-espace vectoriel de dimension finie.  
Un endomorphisme
\[
u\in\mathcal L(E,E)
\]
est dit \textbf{trigonalisable} s’il existe une base de \(E\) dans laquelle la matrice de \(u\) est triangulaire supérieure.
\end{definition}

\subsection{Matrice trigonalisable}

\begin{definition}
Une matrice \(A\in M_n(\K)\) est dite \textbf{trigonalisable} s’il existe \(P\in GL_n(\K)\) et une matrice triangulaire supérieure \(T\) telles que
\[
A=PTP^{-1}.
\]
\end{definition}

\begin{remarque}
La trigonalisation est une réduction moins forte que la diagonalisation :
\[
\text{diagonalisable} \Longrightarrow \text{trigonalisable}.
\]
La réciproque est fausse en général.
\end{remarque}

\section{Sous-espaces stables}

\subsection{Définition}

\begin{definition}
Soit \(u\in\mathcal L(E,E)\).  
Un sous-espace vectoriel \(F\subset E\) est dit \textbf{stable} par \(u\) si
\[
u(F)\subset F.
\]
\end{definition}

\begin{remarque}
Les sous-espaces stables jouent un rôle central dans la trigonalisation, car une matrice triangulaire correspond à l’existence d’un drapeau de sous-espaces stables.
\end{remarque}

\subsection{Drapeau stable}

\begin{definition}
On appelle \textbf{drapeau} de \(E\) une suite de sous-espaces
\[
\{0\}=E_0\subset E_1\subset E_2\subset\cdots\subset E_n=E
\]
telle que
\[
\dim(E_k)=k.
\]
\end{definition}

\begin{theoreme}
Un endomorphisme \(u\) de \(E\) est trigonalisable si et seulement s’il existe un drapeau de \(E\) stable par \(u\).
\end{theoreme}

\begin{proof}
Supposons d’abord \(u\) trigonalisable. Soit \((e_1,\dots,e_n)\) une base dans laquelle la matrice de \(u\) est triangulaire supérieure. Posons
\[
E_k=\Vect(e_1,\dots,e_k).
\]
Comme la matrice est triangulaire supérieure, on a
\[
u(e_k)\in \Vect(e_1,\dots,e_k)=E_k.
\]
Donc chaque \(E_k\) est stable, et on obtient un drapeau stable.

Réciproquement, supposons qu’il existe un drapeau stable
\[
\{0\}=E_0\subset E_1\subset\cdots\subset E_n=E.
\]
Choisissons pour chaque \(k\) un vecteur \(e_k\in E_k\setminus E_{k-1}\). Alors \((e_1,\dots,e_n)\) est une base de \(E\), et comme chaque \(E_k\) est stable,
\[
u(e_k)\in E_k=\Vect(e_1,\dots,e_k).
\]
Donc la matrice de \(u\) dans cette base est triangulaire supérieure.
\end{proof}

\section{Critère fondamental de trigonalisation}

\subsection{Énoncé}

\begin{theoreme}
Un endomorphisme \(u\) d’un \(\K\)-espace vectoriel de dimension finie est trigonalisable sur \(\K\) si et seulement si son polynôme caractéristique est scindé sur \(\K\).
\end{theoreme}

C’est l’un des résultats centraux du chapitre.

\subsection{Sens direct}

\begin{proposition}
Si \(u\) est trigonalisable sur \(\K\), alors son polynôme caractéristique est scindé sur \(\K\).
\end{proposition}

\begin{proof}
Dans une base adaptée, la matrice de \(u\) est triangulaire supérieure :
\[
T=
\begin{pmatrix}
\lambda_1 & * & \cdots & *\\
0 & \lambda_2 & \cdots & *\\
\vdots & \vdots & \ddots & \vdots\\
0 & 0 & \cdots & \lambda_n
\end{pmatrix}.
\]
Alors
\[
\chi_u(X)=\chi_T(X)=\prod_{i=1}^n (X-\lambda_i),
\]
donc le polynôme caractéristique est scindé sur \(\K\).
\end{proof}

\subsection{Sens réciproque}

\begin{theoreme}
Si le polynôme caractéristique de \(u\) est scindé sur \(\K\), alors \(u\) est trigonalisable sur \(\K\).
\end{theoreme}

\begin{proof}
On raisonne par récurrence sur \(n=\dim(E)\).

\textbf{Initialisation :} si \(n=1\), c’est évident.

\textbf{Hérédité :} supposons le résultat vrai jusqu’à la dimension \(n-1\), et soit \(u\in \mathcal L(E,E)\) avec \(\dim(E)=n\), tel que \(\chi_u\) soit scindé.

Comme \(\chi_u\) est scindé, il admet une racine \(\lambda\in\K\). Donc \(\lambda\) est une valeur propre de \(u\), et il existe un vecteur propre non nul \(e_1\) tel que
\[
u(e_1)=\lambda e_1.
\]
Le sous-espace
\[
E_1=\Vect(e_1)
\]
est donc stable par \(u\).

Considérons l’endomorphisme induit \(\widetilde u\) sur le quotient \(E/E_1\). Son polynôme caractéristique est encore scindé, car il correspond au quotient par un facteur \((X-\lambda)\).

Par hypothèse de récurrence, \(\widetilde u\) est trigonalisable. Il existe donc une base de \(E/E_1\) dans laquelle sa matrice est triangulaire supérieure. En relevant cette base dans \(E\), puis en ajoutant \(e_1\), on obtient une base de \(E\) dans laquelle la matrice de \(u\) est triangulaire supérieure.

Ainsi \(u\) est trigonalisable.
\end{proof}

\begin{remarque}
Ce théorème montre que la trigonalisation est exactement la bonne notion quand le polynôme caractéristique est scindé, même en l’absence de diagonalisation.
\end{remarque}

\section{Conséquences importantes}

\subsection{Toute matrice complexe est trigonalisable}

\begin{corollaire}
Toute matrice \(A\in M_n(\C)\) est trigonalisable sur \(\C\).
\end{corollaire}

\begin{proof}
Le polynôme caractéristique de \(A\) est un polynôme complexe de degré \(n\), donc il est scindé sur \(\C\) par le théorème fondamental de l’algèbre. On applique alors le théorème de trigonalisation.
\end{proof}

\begin{remarque}
C’est un résultat fondamental : sur \(\C\), toute matrice peut au moins être triangularisée.
\end{remarque}

\subsection{Sur \(\R\)}

\begin{corollaire}
Une matrice réelle est trigonalisable sur \(\R\) si et seulement si son polynôme caractéristique est scindé sur \(\R\).
\end{corollaire}

\begin{exemple}
La matrice
\[
A=
\begin{pmatrix}
0 & -1\\
1 & 0
\end{pmatrix}
\]
a pour polynôme caractéristique
\[
X^2+1,
\]
qui n’est pas scindé sur \(\R\). Elle n’est donc pas trigonalisable sur \(\R\), mais elle l’est sur \(\C\).
\end{exemple}

\section{Trigonalisation et diagonalisation}

\subsection{Comparaison}

\begin{proposition}
Si un endomorphisme est diagonalisable, alors il est trigonalisable.
\end{proposition}

\begin{proof}
Une matrice diagonale est en particulier triangulaire supérieure.
\end{proof}

\begin{remarque}
La réciproque est fausse : une matrice triangulaire peut ne pas être diagonalisable.
\end{remarque}

\subsection{Exemple fondamental}

\begin{exemple}
La matrice
\[
A=
\begin{pmatrix}
1 & 1\\
0 & 1
\end{pmatrix}
\]
est déjà triangulaire supérieure, donc trigonalisable.  
Mais elle n’est pas diagonalisable, car son unique valeur propre \(1\) a multiplicité algébrique \(2\) et multiplicité géométrique \(1\).
\end{exemple}

\section{Matrice triangulaire par blocs}

\subsection{Principe}

\begin{proposition}
Soit \(u\in\mathcal L(E,E)\), et soit \(F\) un sous-espace stable par \(u\).  
Si l’on choisit une base de \(F\), complétée en une base de \(E\), alors la matrice de \(u\) dans cette base est de la forme
\[
\begin{pmatrix}
A & B\\
0 & C
\end{pmatrix}.
\]
\end{proposition}

\begin{proof}
Comme \(F\) est stable, l’image de tout vecteur de la base de \(F\) reste dans \(F\), donc les coordonnées sur les vecteurs hors de \(F\) sont nulles. Cela force le bloc inférieur gauche à être nul.
\end{proof}

\begin{remarque}
Cette forme triangulaire par blocs est très utile pour les raisonnements par récurrence.
\end{remarque}

\subsection{Polynôme caractéristique}

\begin{proposition}
Si
\[
M=
\begin{pmatrix}
A & B\\
0 & C
\end{pmatrix},
\]
alors
\[
\chi_M(X)=\chi_A(X)\chi_C(X).
\]
\end{proposition}

\begin{proof}
On calcule
\[
XI-M=
\begin{pmatrix}
XI-A & -B\\
0 & XI-C
\end{pmatrix},
\]
et son déterminant est le produit des déterminants des blocs diagonaux, puisque la matrice est triangulaire par blocs.
\end{proof}

\section{Sous-espaces caractéristiques}

\subsection{Première apparition}

\begin{definition}
Soit \(u\in\mathcal L(E,E)\) et \(\lambda\in\K\). Pour \(m\ge 1\), on considère
\[
\Ker\big((u-\lambda\id)^m\big).
\]
Ces espaces jouent un rôle important dans la réduction fine.
\end{definition}

\begin{remarque}
Lorsque \(u\) n’est pas diagonalisable, le sous-espace propre
\[
\Ker(u-\lambda\id)
\]
n’est parfois pas suffisant pour comprendre complètement le comportement de \(u\). On est alors amené à considérer des noyaux de puissances.
\end{remarque}

\begin{remarque}
Dans un premier cours, on peut se contenter d’en garder l’idée sans développer toute la théorie de Jordan.
\end{remarque}

\section{Cas usuels de trigonalisation}

\subsection{Endomorphisme dont le polynôme caractéristique est scindé}

\begin{proposition}
Tout endomorphisme de dimension finie dont le polynôme caractéristique est scindé est trigonalisable.
\end{proposition}

\begin{proof}
C’est exactement le théorème fondamental du chapitre.
\end{proof}

\subsection{Projecteurs, involutions}

\begin{proposition}
Un projecteur est trigonalisable.
\end{proposition}

\begin{proof}
Un projecteur vérifie
\[
u^2=u,
\]
donc son polynôme annulateur
\[
X(X-1)
\]
est scindé. Son polynôme caractéristique est donc scindé, et même l’endomorphisme est diagonalisable.
\end{proof}

\begin{proposition}
Une involution est trigonalisable. Si \(\mathrm{car}(\K)\neq 2\), elle est même diagonalisable.
\end{proposition}

\begin{proof}
Une involution vérifie
\[
u^2=\id,
\]
donc elle est annulée par
\[
X^2-1=(X-1)(X+1).
\]
Le polynôme est scindé, donc \(u\) est trigonalisable. Si les racines sont distinctes, elle est diagonalisable.
\end{proof}

\section{Que gagne-t-on avec une triangulaire ?}

\subsection{Lecture des valeurs propres}

\begin{proposition}
Si \(A\) est trigonalisable, alors ses valeurs propres sont les coefficients diagonaux d’une matrice triangulaire qui lui est semblable.
\end{proposition}

\begin{proof}
Résulte de la propriété des matrices triangulaires.
\end{proof}

\subsection{Trace et déterminant}

\begin{proposition}
Si \(A\) est semblable à une matrice triangulaire
\[
T=
\begin{pmatrix}
\lambda_1 & * & \cdots & *\\
0 & \lambda_2 & \cdots & *\\
\vdots & \vdots & \ddots & \vdots\\
0 & 0 & \cdots & \lambda_n
\end{pmatrix},
\]
alors
\[
\tr(A)=\lambda_1+\cdots+\lambda_n
\]
et
\[
\det(A)=\lambda_1\cdots\lambda_n.
\]
\end{proposition}

\begin{proof}
La trace et le déterminant sont invariants par similitude, et on sait les calculer sur une matrice triangulaire.
\end{proof}

\subsection{Polynômes d’une matrice triangulaire}

\begin{proposition}
Si \(T\) est triangulaire supérieure et \(P\in\K[X]\), alors \(P(T)\) est encore triangulaire supérieure.
\end{proposition}

\begin{proof}
Les puissances de \(T\) sont triangulaires supérieures, et une combinaison linéaire de matrices triangulaires supérieures l’est encore.
\end{proof}

\begin{remarque}
Cela permet souvent de simplifier certains calculs, même sans diagonalisation.
\end{remarque}

\section{Exemples détaillés}

\subsection{Exemple déjà triangulaire}

\begin{exemple}
La matrice
\[
A=
\begin{pmatrix}
2 & 1 & 0\\
0 & 2 & 1\\
0 & 0 & 3
\end{pmatrix}
\]
est déjà triangulaire supérieure.

Son polynôme caractéristique vaut
\[
\chi_A(X)=(X-2)^2(X-3).
\]
Les valeurs propres sont \(2\) et \(3\). La matrice est trigonalisable, ce qui est immédiat, mais n’est pas forcément diagonalisable.

Pour \(\lambda=2\),
\[
A-2I_3=
\begin{pmatrix}
0 & 1 & 0\\
0 & 0 & 1\\
0 & 0 & 1
\end{pmatrix},
\]
et l’on voit que le sous-espace propre associé à \(2\) n’est que de dimension \(1\). Donc \(A\) n’est pas diagonalisable.
\end{exemple}

\subsection{Exemple réelle non trigonalisable sur \(\R\)}

\begin{exemple}
Considérons
\[
A=
\begin{pmatrix}
0 & -1\\
1 & 0
\end{pmatrix}.
\]
On a
\[
\chi_A(X)=X^2+1.
\]
Ce polynôme n’est pas scindé sur \(\R\), donc \(A\) n’est pas trigonalisable sur \(\R\).

En revanche, sur \(\C\),
\[
X^2+1=(X-i)(X+i),
\]
donc \(A\) est trigonalisable sur \(\C\), et même diagonalisable sur \(\C\).
\end{exemple}

\subsection{Exemple complexe}

\begin{exemple}
Toute matrice de \(M_n(\C)\) est trigonalisable.  
Par exemple, même si une matrice complexe n’a pas assez de vecteurs propres pour être diagonalisable, elle reste semblable à une matrice triangulaire supérieure.
\end{exemple}

\section{Aperçu des blocs de Jordan}

\subsection{Idée générale}

\begin{remarque}
Quand une matrice n’est pas diagonalisable, on peut souvent la comprendre comme une matrice presque diagonale, avec des blocs du type
\[
J_r(\lambda)=
\begin{pmatrix}
\lambda & 1 & 0 & \cdots & 0\\
0 & \lambda & 1 & \ddots & \vdots\\
\vdots & \ddots & \ddots & \ddots & 0\\
0 & \cdots & 0 & \lambda & 1\\
0 & \cdots & \cdots & 0 & \lambda
\end{pmatrix}.
\]
\end{remarque}

\begin{remarque}
Dans beaucoup de programmes, on ne développe pas la théorie complète de Jordan à ce stade. Mais ces blocs expliquent la structure des matrices trigonalisables non diagonalisables.
\end{remarque}

\subsection{Bloc de Jordan minimal}

\begin{exemple}
La matrice
\[
\begin{pmatrix}
\lambda & 1\\
0 & \lambda
\end{pmatrix}
\]
est le plus petit exemple d’une matrice trigonalisable non diagonalisable.
\end{exemple}

\section{Critères pratiques}

\begin{theoreme}
Pour une matrice \(A\in M_n(\K)\), les propriétés suivantes sont équivalentes :
\begin{enumerate}[label=\arabic*)]
    \item \(A\) est trigonalisable sur \(\K\) ;
    \item son polynôme caractéristique est scindé sur \(\K\) ;
    \item il existe une base de \(\K^n\) dans laquelle la matrice de \(A\) est triangulaire supérieure.
\end{enumerate}
\end{theoreme}

\begin{proof}
Les équivalences ont été démontrées dans les sections précédentes.
\end{proof}

\section{Méthodes pratiques}

\begin{methode}
Pour savoir si une matrice est trigonalisable sur \(\K\) :
\begin{enumerate}
    \item calculer son polynôme caractéristique ;
    \item vérifier s’il est scindé sur \(\K\).
\end{enumerate}
\end{methode}

\begin{methode}
Pour distinguer diagonalisation et trigonalisation :
\begin{itemize}
    \item si le polynôme caractéristique n’est pas scindé : pas de trigonalisation sur \(\K\) ;
    \item s’il est scindé : trigonalisation possible ;
    \item pour la diagonalisation, il faut en plus suffisamment de vecteurs propres.
\end{itemize}
\end{methode}

\begin{methode}
Pour construire une forme triangulaire :
\begin{enumerate}
    \item chercher une valeur propre \(\lambda\) ;
    \item choisir un vecteur propre \(e_1\) ;
    \item compléter \(e_1\) en une base ;
    \item étudier l’endomorphisme induit sur le quotient ou sur un sous-espace complémentaire ;
    \item procéder par récurrence.
\end{enumerate}
\end{methode}

\section{Erreurs classiques}

\begin{itemize}
    \item Confondre trigonalisable et diagonalisable.
    \item Oublier qu’une matrice réelle peut être trigonalisable sur \(\C\) mais pas sur \(\R\).
    \item Croire qu’une matrice triangulaire est forcément diagonalisable.
    \item Oublier que les valeurs propres d’une matrice triangulaire sont ses coefficients diagonaux.
    \item Penser qu’il faut connaître toute la forme de Jordan pour travailler sur la trigonalisation.
\end{itemize}

\section{Résumé du chapitre}

\begin{itemize}
    \item Une matrice est trigonalisable si elle est semblable à une matrice triangulaire supérieure.
    \item Un endomorphisme est trigonalisable si et seulement s’il existe un drapeau stable.
    \item Une matrice est trigonalisable sur \(\K\) si et seulement si son polynôme caractéristique est scindé sur \(\K\).
    \item Toute matrice complexe est trigonalisable.
    \item Toute matrice diagonalisable est trigonalisable, mais la réciproque est fausse.
    \item La trigonalisation permet déjà de lire les valeurs propres, la trace et le déterminant.
    \item Les matrices non diagonalisables apparaissent naturellement sous forme triangulaire avec des termes non nuls au-dessus de la diagonale.
\end{itemize}

\section*{Exercices d’application immédiate}

\begin{enumerate}[label=\textbf{Exercice \arabic*.}]
    \item Montrer qu’une matrice diagonale est triangulaire supérieure.

    \item Déterminer les valeurs propres de
    \[
    \begin{pmatrix}
    2 & 1 & 0\\
    0 & 2 & 4\\
    0 & 0 & 3
    \end{pmatrix}.
    \]

    \item Étudier si la matrice précédente est diagonalisable.

    \item Montrer qu’une matrice réelle de polynôme caractéristique scindé est trigonalisable sur \(\R\).

    \item Montrer que toute matrice complexe est trigonalisable.

    \item Étudier la trigonalisabilité sur \(\R\) puis sur \(\C\) de
    \[
    \begin{pmatrix}
    0 & -1\\
    1 & 0
    \end{pmatrix}.
    \]

    \item Montrer qu’un projecteur est trigonalisable.

    \item Montrer qu’une involution est trigonalisable.

    \item Donner un exemple de matrice trigonalisable non diagonalisable.

    \item Soit
    \[
    A=
    \begin{pmatrix}
    1 & 1\\
    0 & 1
    \end{pmatrix}.
    \]
    Calculer \(A^n\) pour \(n\in\N\).
\end{enumerate}

\end{document}